\section{Data}
\label{sec:data}

\subsection{Sources}
No public contact-centre logs exist, so we combine two public datasets with real behaviour and a documented
generation step (Table~\ref{tab:sources}).
\begin{itemize}
  \item \textbf{Olist Brazilian e-commerce dataset}~\cite{olist}: 99,441 real, anonymised orders placed
        between 4 September 2016 and 17 October 2018, with items, payments, reviews, customers, sellers,
        products and every delivery timestamp (purchase, approval, carrier hand-off, delivery, estimate).
        Licence CC BY-NC-SA 4.0.
  \item \textbf{Bitext retail e-commerce customer-service dataset}~\cite{bitext}: 44,884 customer messages
        labelled with 46 intents (a customer's goal, such as \code{track\_order}). We keep the
        \textbf{17 intents linked to an order}, grouped into four families: \emph{where is my order}
        (wismo, 4 intents), \emph{cancel/refund} (3), \emph{after delivery} (7) and \emph{order admin} (3).
        Licence CDLA-Sharing 1.0.
  \item \textbf{Knowledge base}: 10 policy documents written for the project (refunds, returns, delivery
        delays, payments, invoices, an advisor playbook...), split into 58 sections.
\end{itemize}

\begin{table}[h]
\centering\small
\caption{Data used by CallVerse (SQLite store and processed tables).}
\label{tab:sources}
\begin{tabularx}{\textwidth}{@{}lrY@{}}
\toprule
\textbf{Table} & \textbf{Rows} & \textbf{Origin and use} \\
\midrule
Orders / customers & 99,441 & Olist; order situation, delivery dates, the advisor's order lookup \\
Order items / payments & 112,650 / 103,886 & Olist; refund amounts, payment anomalies \\
Products / sellers & 32,951 / 3,095 & Olist; product and seller context \\
Reviews & 99,224 & Olist; satisfaction signal during exploration \\
Contacts / tickets & 9,179 / 9,179 & Generated from the orders' situations (Section~\ref{sec:synthetic}) \\
Chat messages & 60,837 & Generated, with the wording of real Bitext messages \\
Knowledge-base sections & 58 & 10 policy documents, embedded with bge-m3 \\
\bottomrule
\end{tabularx}
\end{table}

\subsection{Exploration and quality validation}
Exploration (notebooks 01--02) confirmed that the orders are suitable: 97.0\% are delivered, 1.1\% shipped,
0.6\% cancelled and 0.6\% unavailable. All orders have a purchase time and an estimated delivery date; the
actual delivery time is missing for 2,965 orders, the carrier hand-off for 1,783 and the approval for 160,
which matches their statuses (an undelivered order has no delivery time). Review comments are often empty
(58.7\% have no message), and 610 products (1.9\%) have no category.

The quality validation found five kinds of contradictions (Table~\ref{tab:quality}). One correction matters for
the whole project: the estimated delivery dates are stored at midnight, so comparing them with a delivery
\emph{timestamp} counted orders delivered later on the same day as ``late''. Measured by calendar date,
\textbf{6,535 of 96,476 delivered orders (6.8\%) are late}, not the 7,826 reported in a first version.

\begin{table}[h]
\centering\small
\caption{Data quality findings. Each anomaly becomes a Boolean flag column; no row is deleted or imputed.}
\label{tab:quality}
\begin{tabularx}{0.8\textwidth}{@{}Yr@{}}
\toprule
\textbf{Finding (flag column)} & \textbf{Orders} \\
\midrule
Carrier hand-off before payment approval & 1,359 \\
Delivery before carrier hand-off & 23 \\
Review created before the purchase date (40 of them reuse a review id) & 64 \\
Order with more than one review & 547 \\
Payment total differs from items + freight by more than 1 BRL & 249 \\
\bottomrule
\end{tabularx}
\end{table}

\subsection{Cleaning and feature preparation}
Cleaning (notebooks 03--04) parses types and \emph{flags} anomalies rather than deleting them, so record
counts are preserved and no value is invented by imputation. It produces an \textbf{analytical order table of
77 features} (delivery delay, lateness, payment anomalies, review score, product and seller context...) and a
daily activity table (orders purchased, handed to carriers, delivered, delivered late). Six tables that later
layers depend on are protected by \emph{data contracts}: Pandera schemas~\cite{pandera} check keys, types,
allowed values and ranges (\code{callverse validate}); all pass.

\subsection{Synthetic contacts, tickets and chats}
\label{sec:synthetic}
Real contact logs are private, so contacts are generated from the real orders with explicit, documented
assumptions (\code{configs/data.yaml}, notebooks 05--08 and 12):
\begin{enumerate}
  \item Each order is placed in a \emph{situation}: never delivered and overdue, delivered late, cancelled,
        problem after delivery, payment anomaly, or uneventful.
  \item Each situation has a \emph{contact probability} (for example, 35\% of late deliveries contact
        support), an \emph{intent mix}, and a \emph{timing rule} relative to an order event (a delay-related
        contact comes after the estimated date has passed). Contact hours follow the hour-of-day profile of
        Olist purchases, and a promotion calendar adds known retail events.
  \item The customer's opening message is a real Bitext message of the chosen intent.
\end{enumerate}
The result is \textbf{9,179 contacts} from January 2017 to August 2018. We cross-checked the volume against a
public benchmark (``where is my order'' contacts at 3--10\% of orders). Each contact then becomes a ticket with
an outcome decided by the company rules (resolved in the chat 53.4\%, follow-up needed 43.4\%, escalated
3.2\%), a handle time drawn from the simulation's assumptions, and a satisfaction survey answered 30\% of the
time (mean 3.93/5). The \textbf{60,837 chat messages} use approved answer templates. All generated records carry
an \code{is\_synthetic} flag. Generation is seeded, so it is reproducible.

\subsection{Ethics, licences and confidentiality}
Olist identifiers are already anonymised, and no personal data is collected. Synthetic records are flagged
and the rates behind them are written in one configuration file, so they can be replaced by a client's real
logs without code changes. We respect the licences: non-commercial, share-alike use of Olist (CC BY-NC-SA 4.0)
and Bitext (CDLA-Sharing 1.0). The advisor never asks for card numbers or passwords (a prompt rule and a
checklist item), customers are told they talk to an AI, and staffing recommendations remain human decisions.
For the public deployment, a script removes the customers we created while testing (our own names) and the
transcripts of our test chats before anything leaves the laptop (Section~\ref{sec:deployment}).
